\begin{abstract}
We study how demand-responsive transport supply changes spatial retail price competition. Two retail destinations compete in prices, while a third-party operator reallocates a fixed service resource across directions in response to passenger demand. One direction also carries exogenous background travel demand. A retailer's price therefore affects its market share directly and indirectly by shifting service toward its own direction and away from its rival. We show that a nonempty open set of global pure-strategy Hotelling--Bertrand equilibria exhibits one-sided strategic substitutability: one retailer lowers its price when the rival raises price, while the other retains an upward-sloping response. A minimum service obligation can be strictly slack at equilibrium yet support the global equilibrium by truncating extreme off-path service withdrawal. The operator's waiting-cost minimization also yields an exact welfare mapping between the access term in shopper demand and marginal aggregate waiting cost. The decentralized allocation generally differs from the constrained-efficient allocation under the same service floor. The mechanism is locally robust to power waiting-cost specifications and applies most naturally to managed directional, shuttle, and demand-responsive fleet-allocation settings.
\end{abstract}
